\documentclass{article}
\usepackage{amsmath,amssymb,amsthm}
\usepackage{algorithm}
\usepackage[noend]{algpseudocode}
\usepackage{fullpage}
\usepackage{hyperref}
\newtheorem{theorem}{Theorem}
\newtheorem{lemma}{Lemma}
\newtheorem{assumption}{Assumption}
\newcommand{\gap}{\mathcal{G}}
\newcommand{\diam}{D}
\begin{document}

\section*{Algorithm Name}
\textbf{Frank–Wolfe (Conditional Gradient) for Non‑Convex Smooth Optimization}

\section*{Mathematical Formulation}
Let $f:\mathbb{R}^n\rightarrow\mathbb{R}$ be differentiable and $C\subset\mathbb{R}^n$ a non‑empty compact convex set.
The algorithm generates a sequence $\{x_k\}_{k\ge 0}\subset C$ as follows:

\begin{align}
\text{(Linear minimization oracle)}\qquad
s_k &:= \arg\min_{s\in C}\;\langle \nabla f(x_k),\,s\rangle ,\label{eq:oracle}\\[4pt]
\text{(Update)}\qquad
x_{k+1} &:= (1-\gamma_k)x_k + \gamma_k s_k ,\qquad \gamma_k\in(0,1].\label{eq:update}
\end{align}

The \emph{Frank–Wolfe gap} (a stationarity measure) at $x\in C$ is
\[
\gap(x) \;:=\; \max_{s\in C}\langle \nabla f(x),\,x-s\rangle
               \;=\;\langle \nabla f(x),\,x-s_x\rangle,
\]
where $s_x$ denotes the minimizer in \eqref{eq:oracle}.  Note that $\gap(x)\ge 0$ and $\gap(x)=0$ iff $x$ satisfies the first‑order optimality condition over $C$.

\section*{Pseudocode}
\begin{algorithm}[H]
\caption{Frank–Wolfe for non‑convex $f$}
\begin{algorithmic}[1]
\State \textbf{Input:} initial point $x_0\in C$, step‑size schedule $\{\gamma_k\}_{k\ge0}\subset(0,1]$, maximum iterations $K$.
\For{$k=0,1,\dots,K-1$}
    \State Compute gradient $g_k \gets \nabla f(x_k)$.
    \State \textbf{Linear minimization oracle:}
          \[
          s_k \gets \arg\min_{s\in C}\langle g_k, s\rangle .
          \]
    \State Update $x_{k+1} \gets (1-\gamma_k)x_k + \gamma_k s_k$.
\EndFor
\State \textbf{Output:} $\{x_k\}_{k=0}^K$ (or the iterate with smallest gap).
\end{algorithmic}
\end{algorithm}

\section*{Dry Run Example}
Consider the one‑dimensional problem
\[
f(w) = (w-3)^2,\qquad C=[-5,5],\qquad w_0=0,\qquad \gamma_k\equiv\gamma=0.1.
\]
The gradient is $\nabla f(w)=2(w-3)$ and the diameter $D=10$.

\begin{center}
\begin{tabular}{c|c|c|c|c}
$k$ & $w_k$ & $\nabla f(w_k)$ & $s_k$ (oracle) & $w_{k+1}$\\\hline
0 & $0$ & $-6$ & $\displaystyle\arg\min_{s\in[-5,5]}(-6)s = 5$ & $0.9\cdot0+0.1\cdot5=0.5$\\
1 & $0.5$ & $-5$ & $\displaystyle\arg\min_{s\in[-5,5]}(-5)s = 5$ & $0.9\cdot0.5+0.1\cdot5=0.95$\\
2 & $0.95$ & $-4.1$ & $\displaystyle\arg\min_{s\in[-5,5]}(-4.1)s = 5$ & $0.9\cdot0.95+0.1\cdot5=1.355$\\
\end{tabular}
\end{center}
The objective values are $f(0)=9$, $f(0.5)=6.25$, $f(0.95)=4.2025$, $f(1.355)=2.709\ldots$, showing monotone decrease.

\newpage
\section*{Assumptions}
\begin{assumption}[Smoothness]\label{ass:smooth}
$f$ is $L$–smooth on an open set containing $C$, i.e.
\[
\|\nabla f(x)-\nabla f(y)\|\le L\|x-y\|\quad\forall x,y\in C,
\]
equivalently
\[
f(y)\le f(x)+\langle\nabla f(x),y-x\rangle+\frac{L}{2}\|y-x\|^2\quad\forall x,y\in C.
\]
\end{assumption}

\begin{assumption}[Bounded domain]\label{ass:diam}
$C$ is compact and convex with finite diameter
\[
\diam:=\max_{x,y\in C}\|x-y\|<\infty .
\]
\end{assumption}

\begin{assumption}[Existence of a minimum]\label{ass:opt}
Let $f_*:=\min_{x\in C}f(x)$ and assume the minimum is attained.
\end{assumption}

All subsequent analysis uses only \ref{ass:smooth}–\ref{ass:opt}.

\section*{Main Convergence Theorem}
\begin{theorem}[Convergence rate of Frank–Wolfe for non‑convex smooth $f$]\label{thm:main}
Let $\{x_k\}$ be generated by \eqref{eq:oracle}–\eqref{eq:update} with a \emph{constant} step size $\gamma\in(0,1]$.
Define the Frank–Wolfe gap $\gap(x_k)=\langle\nabla f(x_k),x_k-s_k\rangle$.
Then for any integer $K\ge1$,
\[
\min_{0\le k<K}\gap(x_k)\;\le\;
\frac{f(x_0)-f_*}{\gamma K}\;+\;\frac{L\,\diam^2}{2}\,\gamma .
\tag{1}
\]
Choosing
\[
\gamma^\star:=\sqrt{\frac{2\bigl(f(x_0)-f_*\bigr)}{L\,\diam^2\,K}}
\qquad (\gamma^\star\le 1\text{ for sufficiently large }K)
\]
yields the \emph{optimal} bound
\[
\min_{0\le k<K}\gap(x_k)\;\le\;
\sqrt{\frac{2L\,\diam^2\bigl(f(x_0)-f_*\bigr)}{K}}
\;=\;O\!\left(\frac{1}{\sqrt{K}}\right).
\tag{2}
\]
\end{theorem}

\section*{Theoretical Properties}
\begin{itemize}
\item \textbf{Stationarity guarantee.} $\gap(x_k)=0$ iff $x_k$ satisfies the first‑order optimality condition over $C$; thus the bound (1) quantifies how close the iterates are to a stationary point.
\item \textbf{Hyper‑parameter sensitivity.} The bound is decreasing in $\gamma$ for the first term and increasing for the second term; the optimal $\gamma^\star$ balances these effects, leading to the $O(1/\sqrt{K})$ rate.
\item \textbf{Comparison to baselines.}
  \begin{itemize}
    \item Gradient descent with projection attains $O(1/K)$ for convex $f$ and $O(1/\sqrt{K})$ for non‑convex $f$ under the same smoothness, but requires Euclidean projections.
    \item Frank–Wolfe avoids projections, using only a linear minimization oracle; the price is the same $O(1/\sqrt{K})$ rate in the non‑convex regime.
  \end{itemize}
\item \textbf{Stability.} The iterates remain in $C$ for all $k$ because $C$ is convex and the update is a convex combination of two points in $C$.
\end{itemize}

\section*{Discussion}
The theorem shows that even without convexity the Frank–Wolfe method drives the Frank–Wolfe gap to zero at the same sublinear speed as projected gradient methods.  In practice the linear minimization oracle is often cheaper than a projection, making the algorithm attractive for large‑scale structured domains (e.g., the simplex, trace‑norm ball).

\newpage
\section*{Preliminaries (Definitions and Lemmas)}
\begin{lemma}[Descent Lemma for $L$‑smooth functions]\label{lem:descent}
For any $x,y\in C$,
\[
f(y)\le f(x)+\langle\nabla f(x),y-x\rangle+\frac{L}{2}\|y-x\|^2 .
\]
\end{lemma}
\begin{proof}
Directly from Assumption \ref{ass:smooth}. \hfill $\square$
\end{proof}

\begin{lemma}[Bound on step length]\label{lem:diam}
For any $k$, the displacement satisfies
\[
\|s_k-x_k\|\le \diam .
\]
\end{lemma}
\begin{proof}
Both $s_k$ and $x_k$ belong to $C$, and $\diam$ is the maximal distance between two points of $C$. \hfill $\square$
\end{proof}

\section*{Main Proof (Step‑by‑Step Derivation)}
We prove Theorem \ref{thm:main}.  Throughout the proof we keep the constant step size $\gamma\in(0,1]$.

\begin{enumerate}
\item[Step 1.] Apply Lemma \ref{lem:descent} with $x=x_k$ and $y=x_{k+1}$:
\[
f(x_{k+1})\le f(x_k)+\langle\nabla f(x_k),x_{k+1}-x_k\rangle+\frac{L}{2}\|x_{k+1}-x_k\|^2 .
\tag{3}
\]

\item[Step 2.] Substitute the update \eqref{eq:update}:
\[
x_{k+1}-x_k = \gamma\,(s_k-x_k) .
\tag{4}
\]

\item[Step 3.] Plug \eqref{eq:oracle} into the inner product:
\[
\langle\nabla f(x_k),x_{k+1}-x_k\rangle
   = \gamma\langle\nabla f(x_k),s_k-x_k\rangle
   = -\gamma\,\gap(x_k) .
\tag{5}
\]

\item[Step 4.] Use Lemma \ref{lem:diam} to bound the squared norm:
\[
\|x_{k+1}-x_k\|^2 = \gamma^2\|s_k-x_k\|^2 \le \gamma^2\diam^2 .
\tag{6}
\]

\item[Step 5.] Insert \eqref{eq:5} and \eqref{eq:6} into \eqref{eq:3}:
\[
f(x_{k+1})\le f(x_k)-\gamma\,\gap(x_k)+\frac{L\diam^2}{2}\,\gamma^2 .
\tag{7}
\]

\item[Step 6.] Rearrange \eqref{eq:7} to isolate $\gap(x_k)$:
\[
\gamma\,\gap(x_k)\le f(x_k)-f(x_{k+1})+\frac{L\diam^2}{2}\,\gamma^2 .
\tag{8}
\]

\item[Step 7.] Sum \eqref{eq:8} over $k=0,\dots,K-1$:
\[
\gamma\sum_{k=0}^{K-1}\gap(x_k)
   \le f(x_0)-f(x_K)+\frac{L\diam^2}{2}\,\gamma^2 K .
\tag{9}
\]

\item[Step 8.] Use $f(x_K)\ge f_*$ (Assumption \ref{ass:opt}) to obtain
\[
\gamma\sum_{k=0}^{K-1}\gap(x_k)
   \le f(x_0)-f_*+\frac{L\diam^2}{2}\,\gamma^2 K .
\tag{10}
\]

\item[Step 9.] Divide both sides of \eqref{eq:10} by $\gamma K$:
\[
\frac{1}{K}\sum_{k=0}^{K-1}\gap(x_k)
   \le \frac{f(x_0)-f_*}{\gamma K}
      +\frac{L\diam^2}{2}\,\gamma .
\tag{11}
\]

\item[Step 10.] Since the minimum of a set is bounded by its average,
\[
\min_{0\le k<K}\gap(x_k)\;\le\;\frac{1}{K}\sum_{k=0}^{K-1}\gap(x_k) .
\tag{12}
\]

\item[Step 11.] Combine \eqref{eq:11} and \eqref{eq:12} to obtain the desired bound \eqref{eq:1}:
\[
\boxed{\displaystyle
\min_{0\le k<K}\gap(x_k)\le
\frac{f(x_0)-f_*}{\gamma K}
+\frac{L\diam^2}{2}\,\gamma } .
\tag{13}
\]

\item[Step 12.] Optimize the right‑hand side over $\gamma\in(0,1]$.  The function
\[
\phi(\gamma)=\frac{A}{\gamma K}+B\gamma,\qquad 
A:=f(x_0)-f_*,\; B:=\frac{L\diam^2}{2},
\]
has derivative $-\frac{A}{\gamma^2 K}+B$.
Setting it to zero yields
\[
\gamma^\star=\sqrt{\frac{A}{BK}}=\sqrt{\frac{2\bigl(f(x_0)-f_*\bigr)}{L\diam^2 K}} .
\tag{14}
\]
Plugging $\gamma^\star$ into \eqref{eq:13} gives
\[
\min_{0\le k<K}\gap(x_k)\le
2\sqrt{\frac{AB}{K}}=
\sqrt{\frac{2L\diam^2\bigl(f(x_0)-f_*\bigr)}{K}} .
\tag{15}
\]
This is precisely the bound \eqref{eq:2}.  \hfill $\square$
\end{enumerate}

\section*{Rate Derivation (With Constants)}
From \eqref{eq:15} we can write the explicit constant‑filled rate:
\[
\boxed{
\min_{0\le k<K}\gap(x_k)\;\le\;
\underbrace{\sqrt{2L}_{\text{smoothness}}
\;\underbrace{\diam}_{\text{domain size}}
\;\underbrace{\sqrt{f(x_0)-f_*}}_{\text{initial suboptimality}}}_{\text{combined constant}}
\;\cdot\;K^{-1/2}} .
\]
Thus the Frank–Wolfe gap decreases at the order $K^{-1/2}$.

\section*{Special Cases}
\begin{enumerate}
\item \textbf{Convex $f$.} If $f$ is convex, the Frank–Wolfe gap also upper bounds the primal suboptimality:
\[
f(x_k)-f_* \le \gap(x_k) .
\]
Consequently the same $O(1/\sqrt{K})$ bound applies to the function value, and with the classical step size $\gamma_k=2/(k+2)$ one recovers the $O(1/K)$ rate for convex objectives \cite{Jaggi2013}.

\item \textbf{Line‑search step size.} Choosing $\gamma_k=\arg\min_{\gamma\in[0,1]}f\bigl((1-\gamma)x_k+\gamma s_k\bigr)$ yields a non‑increasing sequence $f(x_k)$ and often improves practical performance, while the worst‑case $O(1/\sqrt{K})$ bound still holds because $\gamma_k\le1$.

\item \textbf{Polytope with bounded diameter $D=2$.} For the probability simplex $\Delta_n$, $D=\sqrt{2}$, and the bound simplifies to
\[
\min_{k}\gap(x_k)\le
\sqrt{L\bigl(f(x_0)-f_*\bigr)}\;K^{-1/2}.
\]
\end{enumerate}

\begin{thebibliography}{9}
\bibitem{Jaggi2013}
M.~Jaggi, \emph{Revisiting Frank–Wolfe: Projection‑free convex optimization}, ICML 2013.
\end{thebibliography}

\end{document}